\documentclass[11pt]{article}

\usepackage[utf8]{inputenc}
\usepackage[spanish]{babel}
\usepackage[margin=2.5cm]{geometry}
\usepackage{amsmath,amssymb}
\usepackage{hyperref}

\title{Turbulencia y modelos de partículas inerciales}
\author{}
\date{}

\begin{document}
\maketitle

\section{Escalas de la turbulencia}

Consideramos un flujo turbulento homogéneo e isótropo, solución de las
ecuaciones de Navier-Stokes incompresibles
\begin{equation}
  \nabla \cdot \mathbf{u} = 0, \qquad
  \frac{D\mathbf{u}}{Dt} \equiv \partial_t \mathbf{u} + (\mathbf{u}\cdot\nabla)\mathbf{u}
  = -\frac{1}{\rho}\nabla p + \nu \nabla^2 \mathbf{u},
\end{equation}
donde $\nu$ es la viscosidad cinemática y $\rho$ la densidad del fluido. La
turbulencia se caracteriza por una amplia separación de escalas: la energía
se inyecta a la escala grande y es transferida, sin disipación apreciable,
a través de un rango inercial de escalas cada vez más pequeñas (la cascada
de Richardson-Kolmogorov) hasta llegar a las escalas donde la viscosidad
finalmente disipa la energía en calor.

\subsection{Escala integral y tiempo de giro grande}

La intensidad de las fluctuaciones de velocidad se mide con la energía
cinética media por unidad de masa
\begin{equation}
  E = \frac{1}{2}\langle v^2\rangle
\end{equation}
y con la velocidad cuadrática media total $U = \langle v^2\rangle^{1/2} =
\sqrt{2E}$ (bajo isotropía, $U = \sqrt{3}\,u_{\rm rms}$ si se prefiere la
velocidad cuadrática media por componente $u_{\rm rms} = \langle
u_1^2\rangle^{1/2}$, usada más abajo para $\mathrm{Re}_\lambda$). La tasa
media de disipación de energía cinética por unidad de masa se obtiene a
partir de la enstrofía:
\begin{equation}
  \varepsilon = 2\nu \langle \omega^2 \rangle.
\end{equation}
La escala integral $L$, la escala espacial más grande donde reside la
mayor parte de la energía (a la cual ésta se inyecta), se calcula a partir
del espectro de energía como el cociente de momentos espectrales
\begin{equation}
  L = \frac{\int (E(k)/k)\,dk}{\int E(k)\,dk}.
\end{equation}
Alternativamente, dado que $L$ es por construcción la escala más grande
del sistema, puede aproximarse directamente por el tamaño de la caja de
simulación, sin necesidad de calcular el espectro.

El tiempo característico de las estructuras de mayor escala (el
\emph{tiempo de giro grande}, o \emph{large-eddy turnover time}) es
\begin{equation}
  \tau_L = \frac{L}{U}.
\end{equation}

\subsection{Rango inercial y microescala de Taylor}

En el rango inercial, $\eta \ll \ell \ll L$, las diferencias de velocidad
entre puntos separados una distancia $\ell$ siguen el escaleo de
Kolmogorov (K41):
\begin{equation}
  u_\ell \sim (\varepsilon \ell)^{1/3},
  \qquad
  \tau_\ell = \frac{\ell}{u_\ell} \sim \varepsilon^{-1/3}\ell^{2/3}.
\end{equation}
La intensidad de la turbulencia suele cuantificarse con el número de
Reynolds basado en la microescala de Taylor $\lambda$,
\begin{equation}
  \lambda = \left(\frac{15\,\nu}{\varepsilon}\right)^{1/2} u_{\rm rms},
  \qquad
  \mathrm{Re}_\lambda = \frac{u_{\rm rms}\,\lambda}{\nu}.
\end{equation}

\subsection{Escalas de Kolmogorov}

En el extremo disipativo del rango inercial, la viscosidad se vuelve
dominante y detiene la cascada. Las escalas de Kolmogorov se construyen
dimensionalmente a partir de $\nu$ y $\varepsilon$, los únicos parámetros
relevantes a esa escala:
\begin{equation}
  \eta = \left(\frac{\nu^3}{\varepsilon}\right)^{1/4}, \qquad
  \tau_\eta = \left(\frac{\nu}{\varepsilon}\right)^{1/2}, \qquad
  u_\eta = (\nu\varepsilon)^{1/4}, \qquad
  k_\eta = \frac{1}{\eta}.
\end{equation}
$\eta$ es la escala espacial por debajo de la cual la disipación viscosa
domina, $\tau_\eta$ es el tiempo de giro asociado a esa escala, $u_\eta$
la velocidad característica correspondiente, y $k_\eta$ el número de onda
de Kolmogorov (nótese que $\eta = u_\eta \tau_\eta$ y que $\tau_\eta$
coincide con la extrapolación de $\tau_\ell$ del rango inercial hasta
$\ell=\eta$).

El tiempo de Kolmogorov $\tau_\eta$ es la escala temporal de referencia
para caracterizar la inercia de una partícula suspendida en el flujo: el
número de Stokes se define como
\begin{equation}
  \mathrm{St} = \frac{\tau_p}{\tau_\eta},
\end{equation}
donde $\tau_p$ es el tiempo de respuesta de la partícula (Sección
\ref{sec:particulas}). $\mathrm{St}\to 0$ corresponde a una partícula que
sigue el flujo como un trazador; $\mathrm{St}$ finito da lugar a
desacoplamiento inercial, muestreo preferencial y clustering.

\section{Modelos de partículas}
\label{sec:particulas}

Cada partícula inercial se caracteriza por su tiempo de respuesta o tiempo
de Stokes $\tau_p = a^2/(3\nu)$ (con $a$ el radio de la partícula), que
mide cuán rápido relaja su velocidad $\mathbf{v}(t)$ hacia la velocidad
del fluido $\mathbf{u}(\mathbf{x}_p,t)$ visto en su posición. Los distintos
modelos de arrastre difieren en qué términos de la ecuación de
Maxey-Riley completa se retienen. A continuación se presentan, en orden
creciente de complejidad, los modelos usados en este proyecto (etiquetas
de \texttt{simuls.yaml} entre paréntesis).

\subsection{Trazador fluido (\texttt{LAG})}

En el límite $\mathrm{St}\to 0$ la partícula no tiene inercia y sigue
exactamente el flujo:
\begin{equation}
  \dot{\mathbf{x}}_p = \mathbf{v}(t) = \mathbf{u}(\mathbf{x}_p,t).
\end{equation}

\subsection{Partícula pesada ingenua / arrastre de Stokes lineal}

El modelo más simple para una partícula pesada retiene únicamente el
arrastre de Stokes lineal (Angriman et al. 2020, ec. 10):
\begin{equation}
  \frac{d\mathbf{x}_p}{dt} = \mathbf{v}(t), \qquad
  \frac{d\mathbf{v}}{dt} = \frac{1}{\tau_p}\left[\mathbf{u}(\mathbf{x}_p,t) - \mathbf{v}\right].
\end{equation}
Es válido para partículas pequeñas y pesadas ($\rho_p \gg \rho_f$), donde
el arrastre domina sobre el gradiente de presión, la masa añadida y la
historia.

\subsection{Maxey-Riley completo (\texttt{MR})}

Cuando la densidad de la partícula es comparable a la del fluido, deben
retenerse además el término de gradiente de presión / masa añadida y el de
flotación. Siguiendo Angriman et al. (2020, ec. 9):
\begin{equation}
  \frac{d\mathbf{v}}{dt} = \frac{1}{T_p}\left[\mathbf{u}(\mathbf{x}_p,t)-\mathbf{v}\right]
  + \frac{3}{2}\frac{\gamma}{1+\gamma/2}\frac{D\mathbf{u}(\mathbf{x}_p,t)}{Dt}
  + \frac{1-\gamma}{1+\gamma/2}\,\mathbf{g},
\end{equation}
con $\gamma = \rho_f/\rho_p$ el cociente de densidades fluido-partícula y
$T_p$ el tiempo de Stokes definido de forma análoga a $\tau_p$. En GHOST
$\gamma$ aparece también dentro del número de Reynolds de la partícula
pero se mantiene siempre en $1$: lo único que se varía entre simulaciones
es el radio $a$ (y por lo tanto $\mathrm{St}$).

Sin el término de historia (Basset-Boussinesq), la ecuación de
Maxey-Riley-Gatignol para una partícula neutralmente flotante se reduce a
(Español et al. 2025, ec. 2.6):
\begin{equation}
  \frac{d\mathbf{v}}{dt} = \frac{1}{\tau_p}\left[\mathbf{u}(\mathbf{x},t)-\mathbf{v}(t)\right]
  + \frac{D\mathbf{u}(\mathbf{x},t)}{Dt}.
  \label{eq:mr-gatignol}
\end{equation}
Esta forma —arrastre lineal más el término de aceleración del fluido
$D\mathbf{u}/Dt$— es la base sobre la que se construyen los modelos de
arrastre no lineal a continuación.

\subsection{Arrastre no lineal (\texttt{NLD}, \texttt{ONLD})}

Para números de Reynolds de partícula $\mathrm{Re}_p$ moderados
($1 \lesssim \mathrm{Re}_p \lesssim 40$), el arrastre de Stokes lineal deja
de ser válido y se reemplaza por una corrección tipo Schiller-Naumann.
Only nonlinear drag (\texttt{ONLD}; Español et al. 2025, ec. 2.7):
\begin{equation}
  \frac{d\mathbf{v}}{dt} = \frac{1+0.15\,\mathrm{Re}_p^{0.687}}{\tau_p}
  \left[\mathbf{u}(\mathbf{x},t)-\mathbf{v}(t)\right],
  \qquad
  \mathrm{Re}_p = \frac{(18\,\tau_p/\nu)^{1/2}\,|\mathbf{u}(\mathbf{x},t)-\mathbf{v}(t)|}{1}.
\end{equation}
\texttt{NLD} agrega a esta misma corrección de arrastre no lineal el
término de aceleración del fluido $D\mathbf{u}/Dt$ de la ecuación
\eqref{eq:mr-gatignol}, es decir, reemplaza el arrastre lineal de esa
ecuación por el arrastre no lineal de arriba. En el segundo set de
simulaciones utilizado en este proyecto \texttt{NLD} y \texttt{ONLD}
coinciden (sin término de aceleración); en el primer set, \texttt{NLD}
incluía ese término extra.

\subsection{Término de historia de Basset-Boussinesq (\texttt{BB})}

La forma más completa de la ecuación de Maxey-Riley agrega, a los
términos anteriores, la fuerza de historia: la respuesta viscosa retardada
del fluido a los cambios pasados de la velocidad relativa
$\mathbf{v}-\mathbf{u}$. Siguiendo la formulación de fuerza hidrodinámica
de Maxey \& Riley (1983) (ver también Bec et al. 2024, ec. 3),
\begin{align}
  \frac{4}{3}\pi a^3 \rho_p \frac{d\mathbf{v}}{dt}
  &= \frac{4}{3}\pi a^3 \rho_f \frac{D\mathbf{u}}{Dt}
  - 6\pi\nu\rho_f a\left[\mathbf{v}-\mathbf{u}(\mathbf{x}_p,t)\right]
  - \frac{4}{3}\pi a^3 (\rho_p-\rho_f)\,\mathbf{g} \nonumber\\
  &\quad - \frac{2}{3}\pi a^3 \rho_f \frac{d}{dt}\left[\mathbf{v}-\mathbf{u}(\mathbf{x}_p,t)\right]
  - 6a^2 \rho_f \sqrt{\pi\nu} \int_{-\infty}^{t}
    \frac{1}{\sqrt{t-s}}\,\frac{d}{ds}\left[\mathbf{v}-\mathbf{u}(\mathbf{x}_p,s)\right] ds.
  \label{eq:bb}
\end{align}
El último término, la integral de Basset-Boussinesq, depende de toda la
historia pasada de la aceleración relativa partícula-fluido, pesada por un
núcleo $\propto (t-s)^{-1/2}$: es una fuerza de memoria, no markoviana, que
distingue a \texttt{BB} de \texttt{MR} (donde este término se descarta,
como se verifica explícitamente en Español et al. 2025 que produce
resultados similares al omitirlo para partículas neutralmente flotantes).

\subsection{Corrección de Faxén (\texttt{FAX})}

Cuando el tamaño de la partícula $a$ no es despreciable frente a la escala
de curvatura del flujo no perturbado, deben agregarse correcciones de
Faxén: términos proporcionales al laplaciano de la velocidad del fluido,
evaluados en el centro de la partícula, que corrigen tanto el arrastre de
Stokes como el término de gradiente de presión. Para el arrastre en
régimen de Stokes estacionario (Maxey \& Riley 1983, ec. 7),
\begin{equation}
  \mathbf{F} = 6\pi a \mu \left\{\mathbf{u}(\mathbf{x}_p,t) - \mathbf{v}(t)\right\}
  + \mu\pi a^3 \left(\nabla^2 \mathbf{u}\right)_{\mathbf{x}_p},
\end{equation}
donde el segundo término es la corrección de Faxén al orden más bajo.
Maxey \& Riley (1983) derivan además la corrección de Faxén análoga para
el flujo de Stokes no estacionario, que debe añadirse consistentemente al
término de masa añadida / gradiente de presión. Estas correcciones son de
orden $(a/\eta)^2$ y por lo tanto solo son relevantes cuando la partícula
deja de ser puntual respecto de la escala de Kolmogorov; \texttt{FAX}
agrega estos términos de curvatura a la ecuación de \texttt{MR}.

\subsection{Resumen}

\begin{center}
\begin{tabular}{ll}
\hline
Etiqueta & Términos retenidos respecto de Maxey-Riley completo \\
\hline
\texttt{LAG}  & ninguno: $\mathbf{v}=\mathbf{u}$ (trazador, $\mathrm{St}=0$) \\
\texttt{MR}   & arrastre de Stokes + masa añadida/gradiente de presión + flotación \\
\texttt{NLD}  & arrastre no lineal + término de aceleración $D\mathbf{u}/Dt$ \\
\texttt{ONLD} & arrastre no lineal solamente \\
\texttt{BB}   & \texttt{MR} + integral de historia de Basset-Boussinesq \\
\texttt{FAX}  & \texttt{MR} + corrección de Faxén (curvatura, $\nabla^2\mathbf{u}$) \\
\hline
\end{tabular}
\end{center}

\end{document}
